\clearpage
\subsection{Примеры работы программы}

На рисунке~\ref{fig:perspective} показана перспективная проекция с заливкой
видимых граней. Оранжевым выделены треугольники, синим — шестиугольники.
На рисунке~\ref{fig:wireframe} включены ортографическая проекция и каркас:
видны все рёбра, включая скрытые. Ориентация объекта и масштаб в обоих случаях
одинаковы. Скриншоты получены из работающего приложения.

\begin{figure}[!ht]
  \centering
  \includegraphics[width=0.88\linewidth,height=0.32\textheight,keepaspectratio]
    {\ReportRoot/assets/screenshots/01-perspective.png}
  \caption{Перспективная проекция с заливкой граней}
  \label{fig:perspective}
\end{figure}

\begin{figure}[!ht]
  \centering
  \includegraphics[width=0.88\linewidth,height=0.32\textheight,keepaspectratio]
    {\ReportRoot/assets/screenshots/02-orthographic-wireframe.png}
  \caption{Ортографическая проекция в каркасном режиме}
  \label{fig:wireframe}
\end{figure}
